% Общая преамбула конспектов теории. Подключается через build/defaults-theory.yaml.
\usepackage{amsmath,amssymb,amsthm,mathtools}

% Плотная вёрстка одной колонкой (двухколоночная проверена и отвергнута: часть
% выкладок шире колонки и печатается за полосу — см. docs/build.md).
\emergencystretch=3em

\theoremstyle{plain}
\newtheorem{theorem}{Теорема}[section]
\newtheorem{lemma}[theorem]{Лемма}
\newtheorem{proposition}[theorem]{Предложение}
\newtheorem{corollary}[theorem]{Следствие}

\theoremstyle{definition}
\newtheorem{definition}[theorem]{Определение}
\newtheorem{example}[theorem]{Пример}
\newtheorem{problem}[theorem]{Задача}

\theoremstyle{remark}
% remark нумеруется в общем счётчике: на замечания ссылаются через \ref, а у
% \newtheorem* счётчика нет — \label внутри такого окружения захватывает
% внешний счётчик, и ссылка молча указывает на чужой объект (проверено: было
% «по замечанию 6.1», совпадающее с номером теоремы 6.1).
\newtheorem{remark}[theorem]{Замечание}
% notation не нумеруется намеренно: на блоки обозначений не ссылаются.
\newtheorem*{notation}{Обозначения}

\renewcommand{\proofname}{Доказательство}

% Обозначения, общие для всех вопросов. Расширять здесь, а не в отдельных
% theory.md, иначе один символ получит два смысла в разных конспектах.
\newcommand{\R}{\mathbb{R}}
\newcommand{\C}{\mathbb{C}}
\newcommand{\N}{\mathbb{N}}
\newcommand{\Z}{\mathbb{Z}}
\newcommand{\Q}{\mathbb{Q}}
\newcommand{\E}{\mathsf{E}}
\newcommand{\Var}{\mathsf{D}}
\newcommand{\Prob}{\mathsf{P}}
\newcommand{\T}{^{\mathsf{T}}}
\newcommand{\norm}[1]{\left\lVert #1 \right\rVert}
\newcommand{\abs}[1]{\left\lvert #1 \right\rvert}
\newcommand{\scal}[2]{\left\langle #1, #2 \right\rangle}
\DeclareMathOperator*{\argmin}{arg\,min}
\DeclareMathOperator*{\argmax}{arg\,max}
\DeclareMathOperator{\sgn}{sign}
\DeclareMathOperator{\rank}{rank}
\DeclareMathOperator{\tr}{tr}
\DeclareMathOperator{\diag}{diag}
\DeclareMathOperator{\supp}{supp}

% Многочлены Чебышёва: T_n — тригонометрическая нормировка cos(n arccos x),
% \Tm_n — приведённая (старший коэффициент 1). Нужны в вопросах 01, 11, 13, 19.
\newcommand{\Tm}{\widetilde{T}}
